\documentclass[10pt,a4paper]{article}
\usepackage{fontspec,xeCJK}
\setmainfont{Times New Roman}
\setCJKmainfont[AutoFakeBold=2,AutoFakeSlant=.2]{SimSun}
\setCJKmonofont{SimSun}
\usepackage[margin=20mm,top=22mm,bottom=22mm]{geometry}
\usepackage{paracol,graphicx,hyperref,xurl,fancyhdr,needspace,xcolor}
\hypersetup{unicode=true,hidelinks}\urlstyle{same}
\setlength{\columnsep}{8mm}\setlength{\parindent}{0pt}\setlength{\parskip}{5pt}
\setlength{\emergencystretch}{3em}\tolerance=2000
\widowpenalty=5000\clubpenalty=5000
\pagestyle{fancy}\fancyhf{}\fancyhead[L]{\small Karak · Coerced Labor and Caste in Bengal}\fancyfoot[C]{\thepage}\setlength{\headheight}{14pt}
\begin{document}

{\Large\bfseries Coerced Labor, Caste, and Aspirations for Liberty in Early Nineteenth-Century Bengal: Revisiting the Antislavery Debate from South Asia\par}\medskip

Anirban Karak. Social Sciences Collegiate Division, The University of Chicago, Chicago, IL, USA. Email: akarak@uchicago.edu. Comparative Studies in Society and History (2026), 1–24. Published online: 22 September 2026.\par
Reading edition based on publisher HTML. Original PDF: 24 pages. All 126 original notes retained after their referring paragraphs. Page labels indicate paragraph starting pages. Original authors © 2026. CC BY 4.0.\par
\url{https://doi.org/10.1017/S001041752610067X}\par
\url{https://creativecommons.org/licenses/by/4.0/}\par
Keywords: slavery; antislavery; South Asia; capitalism; caste; counter-revolution; British Empire; Bengal; agrarian labor; colonial law\par\bigskip
\Needspace{10\baselineskip}\section*{Abstract}
{\scriptsize\color{gray} [p. 1]}\ This article revisits the debate on the relationship between capitalism and abolition from the perspective of South Asia. It argues that bringing the antislavery debate into conversation with the extensive literature on caste opens up new comparative possibilities. Commercial slavery in the early modern Indian Ocean world is now well established, but global histories of slavery have rarely inspired global histories of antislavery. The dominant interpretation of abolition in South Asia remains focused on the contingent requirements of colonial governance. Through an analysis of legal conflicts over the status of slavery in the eastern Indian province of Bengal in the 1820s and 1830s, I argue that this approach obscures the ethical-political significance of abolition in a context shaped by caste and agrarian slavery. The emancipatory potential of abolition, however, was undermined by a colonial state committed to revenue extraction and social order, as well as by the global capitalist context of counter-revolution. By situating the colonial archive within this wider historical conjuncture, the essay shows how both the possibilities and limits of abolition emerged from the contradictions of liberalism and global capitalism. It thereby puts the South Asian case in conversation with broader histories of revolution and counter-revolution.\par
\Needspace{10\baselineskip}\section*{Introduction}
{\scriptsize\color{gray} [p. 1]}\ In the vast literature on capitalism and abolition stretching back to Eric Williams, South Asian social history is conspicuously absent.\nobreak\textsuperscript{1} This can partly be explained by the fact that the study of slavery and capitalism was an Atlantic-centered field for much of the twentieth century. In recent decades, however, commercial slavery in the Indian Ocean world has become a prominent field of inquiry.\nobreak\textsuperscript{2} New Histories of Capitalism and studies of “racial capitalism” have contributed further to a widening of our spatial visions and provoked debates about the global, imperial, and racial dimensions of modern capitalism.\nobreak\textsuperscript{3}\par{\small 1. Eric Williams, \textit{Capitalism and Slavery} (University of North Carolina Press, 1944); David Brion Davis, \textit{The Problem of Slavery in Western Culture} (Cornell University Press, 1966); David Brion Davis, \textit{The Problem of Slavery in the Age of Revolution, 1770}–\textit{1823} (Cornell University Press, 1975); Thomas L. Haskell, “Capitalism and the Origins of the Humanitarian Sensibility, Part I,” \textit{The American Historical Review} 90, 2 (1985): 339–61; Thomas L. Haskell, “Capitalism and the Origins of the Humanitarian Sensibility, Part II,” \textit{The American Historical Review} 90, 3 (1985): 547–66; Thomas Bender, ed., \textit{The Antislavery Debate: Capitalism and Abolitionism as a Problem in Historical Interpretation} (University of California Press, 1992).\par}{\small 2. Kate Ekama, Lisa Hellman, and Matthias van Rossum, eds., \textit{Slavery and Bondage in Asia, 1550–1850: Towards a Global History of Coerced Labor} (De Gruyter, 2022); Alessandro Stanziani, \textit{Bondage: Labor and Rights in Eurasia from the Sixteenth to the Early Twentieth Centuries} (Berghahn Books, 2014); Titas Chakraborty and Matthias van Rossum, “Slave Trade and Slavery in Asia—New Perspectives,” \textit{Journal of Social History} 54, 1 (2020): 1–14; Matthias van Rossum, “Slavery and its Transformations: Prolegomena for a Global and Comparative Research Agenda,” \textit{Comparative Studies in Society and History} 63, 3 (2021): 566–98.\par}{\small 3. Sven Beckert, \textit{Empire of Cotton: A Global History} (Knopf, 2014); Edward E. Baptist, \textit{The Half Has Never Been Told: Slavery and the Making of American Capitalism} (Basic Books, 2014); Sven Beckert and Seth Rockman, eds., \textit{Slavery’s Capitalism: A New History of American Economic Development} (University of Pennsylvania Press, 2016); Destin Jenkins and Justin Leroy, eds., \textit{Histories of Racial Capitalism} (Columbia University Press, 2021); Diana Paton, “Gender History, Global History, and Atlantic Slavery: On Racial Capitalism and Social Reproduction,” \textit{The American Historical Review} 127, 2 (2022): 726–54.\par}
{\scriptsize\color{gray} [p. 2]}\ Nevertheless, the new literature on global histories of slavery has not inspired attempts to write global histories of antislavery. Some scholars have called for a better understanding of the slow “drift in community values” that drives authorities to redefine what constitutes coercion, and the need for a global approach to antislavery has occasionally been acknowledged.\nobreak\textsuperscript{4} Yet, the default interpretation of antislavery in South Asia is one that reads it entirely as a function of the contingent requirements of colonial governance.\nobreak\textsuperscript{5} Such an approach ignores the possibility that abolition in colonial contexts, despite its limits, had a genuine ethical purchase on the lives of those it affected. Agrarian social relations in the subcontinent were grounded in practices of caste, which involved differential access to markets, land, and water, the control of (bonded) labor, and the regulation of sexuality. As a modality of power, caste was “sited at a depth within the indigenous society, well beyond the reach of the disciplinary arm of the colonial state.”\nobreak\textsuperscript{6} This dimension has often been concealed, as this essay will also show, through a blending of “indigenous feudal ideology…with colonialist anthropology.”\nobreak\textsuperscript{7}\par{\small 4. David Eltis, Stanley L. Engerman, Seymour Drescher, and David Richardson, “Introduction,” in David Eltis \textit{et al.}, eds., \textit{The Cambridge World History of Slavery, Volume 4, AD 1804–AD 2016} (Cambridge University Press, 2017), 3–19; Seymour Drescher, “Approaches to Global Antislavery,” in David Stefan Doddington and Enrico Del Lago, eds., \textit{Writing the History of Slavery} (Bloomsbury Academic, 2022), 113–32; William Mulligan, “Injection: The Global Spread of Abolitionism,” in Damian A. Pargas and Juliane Schiel, eds., \textit{The Palgrave Handbook of Global Slavery Throughout History} (Palgrave Macmillan, 2023), 553–64.\par}{\small 5. Andrea Major, \textit{Slavery, Abolitionism, and Empire in India, 1772–1843} (Liverpool University Press, 2012); Andrea Major, \textit{Reimagining Empire in India: George Thompson, Anti-Slavery Activism, and the Global Networks of British Colonial Reform, 1831–1858} (Bloomsbury, 2025); Indrani Chatterjee, “Abolition by Denial: The South Asian Example,” in Gwyn Campbell, ed., \textit{Abolition and its Aftermath in Indian Ocean, Africa, and Asia} (Routledge, 2005), 137–153; Indrani Chatterjee, “British Abolitionism from the Vantage of Pre-Colonial South Asian Regimes,” in David Eltis \textit{et al}., eds.\textit{, The Cambridge World History of Slavery, Volume 4, AD 1804–AD 2016} (Cambridge University Press, 2017), 441–65.\par}{\small 6. Ranajit Guha, “Chandra’s Death,” in Ranajit Guha, ed., \textit{Subaltern Studies V} (Oxford University Press, 1987), 135–165, 155.\par}{\small 7. Ibid., 144.\par}
{\scriptsize\color{gray} [p. 3]}\ Hence, evacuating abolition in India of its ideological content \textit{in toto} produces a major historiographical blind spot.\nobreak\textsuperscript{8} It occludes from vision longer-running South Asian social histories and the ethical-political valence of abolition in a context shaped by caste and agrarian slavery. Building on Thomas Holt’s insightful distinction between abolitionism as a social ideology irreducible to class interest, and abolitionist “outcomes” produced by elite-imperial policy, the following sections gradually unpack my key argument and its comparative implications.\nobreak\textsuperscript{9} The next section situates the essay within the broader historiography of capitalism and abolition, and motivates the need to incorporate the specificities of caste and slavery in South Asia into this literature. Section III presents the evidence: a series of legal conflicts and the testimonies of alleged “slaves” in the eastern Indian province of Bengal in the 1820s and 1830s. These disputes show, I argue in Section IV, that slavery was neither a matter of “kinship” nor “gentle” and harmonious.\nobreak\textsuperscript{10} Agrarian slavery existed in both Mughal South Asia and during the early British colonial era (1765–1858), and there were political struggles over the legal status of slavery in the early nineteenth century. Such aspirations for liberty cannot be explained by colonial state-building alone. Acts of defiance, which included claims of having been “wage” workers rather than “slaves,” and an insistence on the right of property to homesteads, were challenges to the practical basis of caste. Though not textually articulated as such, reading the actions of alleged slaves with reference to a larger structure of inference enables us to better grasp the emergence of an anti-caste and eventually Dalit political subjectivity in South Asia from the 1860s onward.\nobreak\textsuperscript{11}\par{\small 8. I use the term “ideology” to refer to a particular, systematic conjuncture of shared ideas, presuppositions, and sentiments that mediates between our objective experience and our subjective interpretation of that experience. In addition to the work of Thomas Holt cited below, I find Barbara Fields’s historical approach to the ideology of racism particularly useful. See Barbara Fields, “Slavery, Race, and Ideology in the United States of America,” \textit{New Left Review} I/181 (1990): 95–118.\par}{\small 9. Thomas C. Holt, \textit{The Problem of Freedom: Race, Labor, and Politics in Jamaica and Britain, 1832–1938} (Johns Hopkins University Press, 1992), 3–53, \textit{passim.}\par}{\small 10. Colonial officials regularly used the trope of “gentle slavery” to deny the need for legal abolition in India. More recently, Indrani Chatterjee, \textit{Gender, Slavery, and Law in Colonial India} (Oxford University Press, 1999) has argued that there was a “slavery-to-kinship continuum” in early colonial India. Both claims are discussed and criticized in greater detail in section IV.\par}{\small 11. “Dalit” means “ground down” or “broken.” The Dalit leader B. R. Ambedkar first used the term to refer to the “depressed classes” and untouchables of India in 1928, but there is a longer-run story to be told about the coming-into-being of an anti-caste and eventually Dalit political consciousness from the mid-nineteenth century onward. Representative works charting various aspects of this history include Rosalind O’Hanlon, \textit{Caste, Conflict, and Ideology: Mahatma Jotirao Phule and Low Caste Protest in Nineteenth-Century Western India} (Cambridge University Press, 1985); Rupa Viswanath, \textit{The Pariah Problem: Caste, Religion, and the Social in Modern India} (Columbia University Press, 2014); Anupama Rao, \textit{The Caste Question: Dalits and the Politics of Modern India} (University of California Press, 2009); Sekhar Bandyopadhyay, \textit{Caste, Protest, and Identity in Colonial India: The Namasudras of Bengal}, \textit{1872–1947} (Oxford University Press, 2011); and Ramnarayan S. Rawat, \textit{Reconsidering Untouchability: Chamars and Dalit History in North India} (Indiana University Press, 2011).\par}
{\scriptsize\color{gray} [p. 3]}\ The ethical-political potential of abolition, however, was undermined by abolitionist policies enacted by a colonial state deeply invested in “order” and revenue extraction. Indeed, Act V of 1843, which put an end to the legal recognition of slavery but did not actively emancipate the enslaved, can scarcely be called “abolition.” Many historians have recognized and emphasized this point, but the larger, global capitalist context of counter-revolution has not received its due. When “abolition” was effected in India, the forces of reaction were ascendant throughout the wider imperial world \textit{and} within Europe and Britain.\nobreak\textsuperscript{12} Colonial officials who governed India also opposed democratic, working-class movements in Britain. “Reform” in the metropole was an attempt to secure the influence of the landed elite in Parliament and to rein in more radical possibilities.\nobreak\textsuperscript{13} Recognizing the need to situate colonial Jamaica within this broader context, Thomas Holt argued that the gradual rise of racist discourses must be historicized with reference to a context where mass political liberty and the control of labor had become a problem for an entrenched (hence counter-revolutionary) bourgeois-capitalist class.\nobreak\textsuperscript{14}\par{\small 12. “Counter-revolution” is, of course, Marx’s term, and an earlier generation of historians regularly used the concept to contextualize the failures of abolition in the Atlantic world. See Karl Marx, “The Bourgeoisie and the Counter-Revolution I–IV,” in David Fernbach, ed., \textit{The Revolutions of 1848: Political Writings Volume 1} (Verso, 1973), 186–212; C. L. R. James, \textit{The Black Jacobins: Toussaint L’Ouverture and the San Domingo Revolution} (Vintage Books, 1963); W. E. B. Du Bois, \textit{Black Reconstruction: An Essay Toward a History of the Part Which Black Folk Played in the Attempt to Reconstruct Democracy in America, 1860–1880} (Harcourt, Brace, and Co., 1935).\par}{\small 13. E. P. Thompson, \textit{The Making of the English Working Class} (Vintage Books, 1963); Boyd Hilton, \textit{A Mad, Bad, and Dangerous People? England, 1783–1846} (Oxford University Press, 2006).\par}{\small 14. Holt, \textit{The Problem of Freedom.}\par}
{\scriptsize\color{gray} [p. 4]}\ I make a similar argument for South Asia by historicizing the wider context of abolition and the colonial archive that is my evidentiary basis. I neither read the archive uncritically nor reject it \textit{tout court.} Instead, the essay grounds both the possibilities and limits of abolition in the contradictions of liberalism and global capitalism, and asks where South Asia fits within broader histories of capitalism, revolution, and counter-revolution.\par
\Needspace{10\baselineskip}\section*{Abolition, Capitalism, and the Question of Ideology}
{\scriptsize\color{gray} [p. 4]}\ The debates over antislavery in the Atlantic world were intense and impassioned, but all parties agreed on a fundamental point: abolitionism was a major ideological shift, and it was one, if not \textit{the} most significant, change that set apart the Age of Revolution from preceding epochs. The animating question of the debate was simple yet sweeping: given the remarkable acceptance and stability of slavery across several millennia, how can we explain the emergence of abolitionism and its political success in the late eighteenth and early nineteenth centuries? No one denied that capitalism had something to do with it, but disagreements abounded about how to conceptualize the link and how to judge the historical and ethical legacy of abolition.\par
{\scriptsize\color{gray} [p. 4]}\ Eric Williams proposed an elegant, two-part thesis: Atlantic slavery financed British industrialization and was thus foundational for British capitalism; abolition occurred when slavery was no longer profitable and capitalists pushed for “free” wage labor.\nobreak\textsuperscript{15} Both halves of the thesis were challenged. Economic historians showed that profits from slavery were far less important for British industrialization than Williams assumed, and mechanization was driven by other factors including domestic wage levels and contingent scientific advances.\nobreak\textsuperscript{16} Seymour Drescher, on the other hand, demonstrated that abolition occurred at a time when profits from slavery were at an all-time high.\nobreak\textsuperscript{17} Abolition, in short, was “econocide.”\par{\small 15. Williams, \textit{Capitalism and Slavery.}\par}{\small 16. Stanley L. Engerman, “The Slave Trade and British Capital Formation in the Eighteenth Century: A Comment on the Williams Thesis,” \textit{The Business History Review} 46, 4 (1972): 430–43; David Eltis and Stanley L. Engerman, “The Importance of Slavery and the Slave Trade to Industrializing Britain,” \textit{The Journal of Economic History} 60, 1 (2000): 123–44. Cf. Kenneth Morgan, \textit{Slavery, Atlantic Trade and the British Economy, 1660–1800} (Cambridge University Press, 2001).\par}{\small 17. Seymour Drescher, \textit{Econocide: British Slavery in the Era of Abolition} (University of North Carolina Press, 1977).\par}
{\scriptsize\color{gray} [p. 5]}\ Some years later, Thomas Haskell re-opened the question of the link between capitalism and humanitarianism in an extended, two-part essay.\nobreak\textsuperscript{18} He argued that abolition became a practical possibility because the expansion of abstract, formally free market relations enlarged people’s capacity to recognize extended chains of moral responsibility. Thomas Bender rightly observed, however, that it was unclear exactly what Haskell meant by the “market,” and whether he was claiming a new way of valu\textit{ing}, or new \textit{values.} \nobreak\textsuperscript{19} The link between capitalism and humanitarianism remained thin and abstract, and the contingent, political trajectory of abolitionism was left underexamined. Haskell pleaded that he was making a case for “enabling conditions” rather than inevitability, but one could be forgiven for thinking that he was smuggling in economic determinism through the back door.\par{\small 18. Haskell, “Capitalism and the Origins of the Humanitarian Sensibility, Parts I and II.”\par}{\small 19. Bender, “Introduction,” in \textit{The Antislavery Debate}, 11, 7.\par}
{\scriptsize\color{gray} [p. 5]}\ By contrast, David Brion Davis, in his two magisterial works on slavery and abolition, emphasized both the newness of abolitionism and its moral-political force as an \textit{ideology.} \nobreak\textsuperscript{20} He showed that a fundamental component of the experience and world-view of the emerging capitalist class that led the abolitionist movement was the question of how to define “freedom” in a rapidly changing world. “Slavery” was a clear-cut and concrete foil, which could be used to symbolize “all the forces that threatened the true destiny of man.”\nobreak\textsuperscript{21} The ideology of freedom and a principled antislavery position, therefore, could not be reduced to the economic interests of capitalists, but nor could the political trajectory of abolitionism be understood without reference to the actions of elite figures.\par{\small 20. Davis, \textit{The Problem of Slavery in Western Culture}; Davis, \textit{The Problem of Slavery in the Age of Revolution, 1770–1823.}\par}{\small 21. Davis, \textit{The Problem of Slavery in the Age of Revolution}, 41.\par}
{\scriptsize\color{gray} [p. 5]}\ Among the early contributors to the debate, Seymour Drescher’s and Thomas Holt’s works offer further insights into how the gap between enabling discourses and actual outcomes might be bridged. Drescher showed that abolitionism became a mass social movement in Britain because it was widely supported by the working classes.\nobreak\textsuperscript{22} Their empathy for the enslaved was rooted in experiences of exploitation as proletarian workers: they used antislavery ideals to resist further proletarianization and not to uncritically support capitalist expansion. Stadial notions of capitalist development rooted in methodological nationalism, therefore, are ill-equipped to grasp the historical and political complexity of abolitionism, which could have found relevance in purportedly “advanced” as well as “backward” regions.\nobreak\textsuperscript{23}\par{\small 22. Seymour Drescher, \textit{Capitalism and Antislavery: British Mobilization in Comparative Perspective} (Oxford University Press, 1987), 1–24, 72–73, 130–32.\par}{\small 23. A recent critique of “methodological nationalism” from a South Asian perspective is Manu Goswami, \textit{Producing India: From Colonial Economy to National Space} (University of Chicago Press, 2004).\par}
{\scriptsize\color{gray} [p. 6]}\ Holt sharpened these insights by distinguishing between abolitionism as a social movement and ideology, and abolitionist outcomes shaped by elite politicians and imperial administrators. This enabled him to argue that antislavery was a revolutionary ideology, and even an abolition tempered by a period of “apprenticeship” had a positive impact on the lives of those emancipated. Nevertheless, imperial policy driven by elite interests was always responding to changed contexts. At a historical moment when mass political liberty had become a problem for the entrenched bourgeois-capitalist class in both the metropole and the colonies, and when the need to motivate “lazy” workers had become a recurring theme, emancipated laborers in Jamaica could not be allowed to become independent freeholders exercising franchise. Hence, biologistic theories of “race” became the alibi for repression and disenfranchisement.\par
{\scriptsize\color{gray} [p. 6]}\ Davis, Drescher, and Holt provide a comparative framework for studying abolition beyond the Atlantic world. Yet, the antislavery debate has rarely been invoked when thinking about questions of caste, slavery, and abolition in South Asia. One reason for this is that the contemporary literature on capitalism, slavery, and empire views “capitalism” in primarily negative terms, as historically tied \textit{only} and often \textit{necessarily} to slavery, extractive empires, “racialization,” and new forms of gendered domination. This tendency is discernible across the board in the new histories of capitalism, in the large literature on global slavery beyond the Atlantic world, and in the much smaller literature on global abolitionism.\nobreak\textsuperscript{24} The dialectical approach of an earlier era has been jettisoned in favor of one-sided formulations that threaten to erase both the newness and the ideological force of abolitionism in the Age of Revolution.\par{\small 24. Representative titles include Walter Johnson, \textit{River of Dark Dreams: Slavery and Empire in the Cotton Kingdom} (Harvard University Press, 2013); Jenkins and Leroy, \textit{Histories of Racial Capitalism}; Paton, “Gender History, Global History, and Atlantic Slavery”; Marcus Rediker, Titas Chakraborty, and Matthias van Rossum, eds., \textit{A Global History of Runaways: Workers, Mobility, and Capitalism}, \textit{1600–1850} (University of California Press, 2019); and Sheetal Chhabria, “Where Does Caste Fit in a Global History of Racial Capitalism?” \textit{Historical Materialism} 31, 2 (2023): 136–60. For critiques of claims about slavery and racism being \textit{necessary} for capitalism, see Trevor Burnard and Giorgio Riello, “Slavery and the New History of Capitalism,” \textit{Journal of Global History} 15, 2 (2020): 215–44; and Loïc Wacquant, “The Trap of ‘Racial Capitalism,’” \textit{European Journal of Sociology} 64, 2 (2023): 153–62.\par}
{\scriptsize\color{gray} [p. 6]}\ The existence of widespread commercial slavery in the early modern Indian Ocean world is now universally acknowledged, but we are only beginning to understand how caste was implicated in this process. Some scholars have shown that, in practice, caste and slavery were interrelated, but the social history of caste beyond its colonial framing has been left underexamined.\nobreak\textsuperscript{25} Moreover, “resistance” is conceptualized far too often as challenges only to extractive European empires.\nobreak\textsuperscript{26} By leaving open the question of whether resistance was also directed at contradictions within indigenous society, and, if so, with what content, historians have left open the possibility of an uncritical affirmation of pre-colonial South Asian societies.\par{\small 25. Nira Wickramasinghe and Alicia Schrikker, “The Ambivalence of Freedom: Slaves in Jaffna, Sri Lanka, in the Eighteenth and Nineteenth Centuries,” \textit{The Journal of Asian Studies} 78, 3 (2019): 497–519.\par}{\small 26. Kate Ekama, “Precarious Freedom: Manumission in Eighteenth-Century Colombo,” \textit{Journal of Social History} 54, 1 (2020): 88–108; Titas Chakraborty, “Slave Trading and Slave Resistance in the Indian Ocean World: The Case of Early Eighteenth-Century Bengal,” \textit{Slavery and Abolition} 40, 4 (2019): 706–26; Sujit Sivasundaram, \textit{Waves Across the South: A New History of Revolution and Empire} (University of Chicago Press, 2020).\par}
{\scriptsize\color{gray} [p. 7]}\ One example can illustrate the stakes. In the now-established and vibrant field of Dalit studies, scholars often emphasize that Dalits found “colonial” liberalism and modernity useful for addressing their own political concerns, including the problem of representation.\nobreak\textsuperscript{27} It is also true, however, that the colonial state always allied with upper-caste landed elites to extract revenue and control Dalit labor in the nineteenth century, and it supported Dalit political claims in the early twentieth century solely to undermine nationalist claims to unity.\nobreak\textsuperscript{28} Moreover, by the 1910s, liberal arguments about property could easily be mobilized to justify the status quo, that is, a caste-differentiated distribution of property.\nobreak\textsuperscript{29} Therefore, to avoid the extremes of \textit{either} a one-sided, ahistorical defense of liberalism \textit{or} an equally one-sided rejection of it \textit{tout court}, we must historicize both the potentials and the failures of liberty in South Asia. By combining insights from the Atlantic-centered antislavery debates with the extensive literature on caste, slavery, and agrarian social relations within South Asia, the next two sections sketch the possible outlines for such a historicization.\par{\small 27. See the introductory chapter of Ramnarayan S. Rawat and K. Satyanarayana, eds., \textit{Dalit Studies} (Duke University Press, 2016), 1–30, 16.\par}{\small 28. Viswanath, \textit{The Pariah Problem}; Sen, \textit{The Decline of the Caste Question.}\par}{\small 29. Rao, \textit{The Caste Question}, chapter 2.\par}
\Needspace{10\baselineskip}\section*{The Evidence: Aspirations for Liberty and the Struggle over the Legal Status of Slavery in Bengal, ca. 1813–1843}
{\scriptsize\color{gray} [p. 7]}\ On 9 September 1826, in the civil court of Dhaka in eastern Bengal, Golak Narayan Ray filed a suit against Kewal Ram Deo and sixteen others. Ray claimed that Kewal Ram and the others were the descendants of Dakai, Puchai, and Manai, the “hereditary slaves” of Ray’s ancestors.\nobreak\textsuperscript{30} According to Ray, the defendants and their ancestors had rendered service as slaves to him and his ancestors for generations, “being supported by lands assigned.” The defendants had recently left Ray’s estate and relocated elsewhere. Ray argued that by local custom they remained his hereditary slaves and sued to recover them. The court accepted the suit, and a summons was issued to the defendants.\par{\small 30. The entire case history can be found in the Indian Law Commission, \textit{Report on Slavery in India with Appendices, Volume 2} (India: 1841): 251–55.\par}
{\scriptsize\color{gray} [p. 7]}\ Appearing for the defendants, Jagannath and Sambunath claimed that their ancestors were rent-paying tenants in the landed estate of Ray and that their “profession” was “service.” Their father was an agent in the estate of the Ray family, and they had never held \textit{nankar} lands for sustenance.\nobreak\textsuperscript{31} At present, the defendants paid Ray, as their “superior landlord,” the “yearly rent of 358 Rupees and 9 annas.” As such, they had full property rights over their person and labor, and Ray’s claim about desertion did not stand. The cause of the complaint, the defendants insisted, was their refusal to hand over to Ray a plot of land that the landlord wished to include in his garden, but in which the defendants had a dwelling of their own. Out of resentment, Ray had brought legal action.\par{\small 31. The word \textit{nankar} is derived from the Persian words \textit{nān} (bread) and \textit{kar} (tax or service). Lands given to slaves, peasants, and laborers in return for regular service were called \textit{nankar} lands.\par}
{\scriptsize\color{gray} [p. 7]}\ Two other defendants, Kewal Ram and Deb Chand, proffered similar defenses, claiming to be tenants rather than slaves. On 23 May 1828, the case came to trial before Mr. D. B. Morrieson. On this occasion, Golak Ray produced an \textit{ikrar} (pledge or promise) from 1790, which had been executed by Dakai, Puchai, and Manai, and which bore the signature of Mr. W. Douglas, the collector of Jalalpur at the time. Based on this evidence and the testimony of witnesses, Mr. Morrieson passed judgment in favor of the plaintiff and ordered the defendants to be made over to Ray as his slaves. In his judgment, Morrieson pointed out that in the ikrar of 1790, Dakai and his brothers had admitted that they were “hereditary slaves” of the plaintiff’s family. It also showed that Ray’s grandfather had “bought the taluka (tenancy) in their name because they were slaves,” that he “fixed the yearly rent at 370 Rupees 2 annas, and that after allowing them nine rupees from the established rent assets as their nankar for services as slaves, he made it over to their charge.” Further, the ikrar contained a clause that Dakai and his brothers \textit{and} their descendants would “continue to render service of slaves to the plaintiff,” and that in case of default, the plaintiff had the right to sue.\par
{\scriptsize\color{gray} [p. 8]}\ The defendants then filed an appeal at the Provincial Court of Appeal, and their case was heard on 9 November 1829, by Mr. W. Cracoft. The judge reversed the decree of the lower court, and he even cast doubt on the ikrar of 1790 that had been the key evidentiary basis of the previous judgment. Cracoft found the deed suspicious, doubted its provenance, and concluded that the defendants were more likely tenants than slaves. His position, however, was contradicted by Mr. Charles Smith, a judge of the same court, who next heard the case on 29 December 1829.\par
{\scriptsize\color{gray} [p. 8]}\ Smith believed the ikrar to be true, and he saw no reason to question its veracity. His sweeping comments demonstrate the presumptive attitudes of many judges: “Slavery of a family may be \textit{inferred from continuous service}; and it seldom happens, that after the lapse of many years, the original title, shewing acquisition of the slave’s forefather, is forthcoming. In this part of the country, many slaves are apparently persons of respectability and educated, and manage the \textit{zemindari} [estate] affairs of their masters. But this constitutes no ground for emancipation. In a word, the appellants and their ancestors are the hereditary slaves of the respondent, and if discharged, notwithstanding proof of their servile state, most slaves will become recusant, and on various pretexts will find means to effect their discharge.”\nobreak\textsuperscript{32}\par{\small 32. Indian Law Commission, \textit{Report on Slavery, Volume 2}, 254, emphasis mine.\par}
{\scriptsize\color{gray} [p. 8]}\ Due to such disagreements among judges, the case went to the Murshidabad Court of Appeal to be heard by a third judge. On 20 April 1830, Mr. C. W. Steer confirmed the judgment of Mr. Smith, and the appellants then requested the highest civil court—the Sadar Dewani Adalat (SDA)—for permission to file a special appeal. This was granted on 21 June 1830, by Mr. Alexander Ross and Mr. R. H. Rattray.\par
{\scriptsize\color{gray} [p. 8]}\ The tone of the SDA judges was diametrically opposed to many in the lower courts. Mr. Rattray, on 26 March 1832, concurred with Mr. Cracoft and with the grounds of his judgment. Moreover, when the lawyers of Ray pointed out that one of the appellants, Kalikinkar, had died, and insisted that his heirs be summoned to follow up on the appeal, Mr. Rattray argued that there was no need to summon them if “death had emancipated [Kalikinkar].” Mr. Ross, on 5 May 1832, finalized the judgment of Mr. Rattray and added that the “service” of slaves and of dependent tenants had to be clearly distinguished. He felt that the appellants had been the latter, who held land “on condition of paying a fixed rent and rendering service.” As long as they held the lands and gained sustenance from them, they had to render service, but “not after abandoning the tenure and thereby discharging themselves.” Thus, “he who has the power to emancipate himself, cannot be considered a slave.”\par
{\scriptsize\color{gray} [p. 9]}\ Such conflicts were extremely common. While peasant rebellions and the question of peasant consciousness in early colonial Bengal have attracted scholarly attention for decades, the actions of slaves and the nature of their claims in legal cases have not been systematically studied.\nobreak\textsuperscript{33} Only a few monographs exist on slavery in early colonial Bengal, and none of them address the question of slave consciousness directly.\nobreak\textsuperscript{34} Yet, alleged “slaves” often expressed outrage against their masters. In the year 1813 alone, they filed 150 cases in colonial courts, which may have represented only the tip of the iceberg.\nobreak\textsuperscript{35} A. D. Campbell, who served as a Superintendent of Police and then as a Magistrate, noted that a slave’s “peculiar situation, want of relatives and friends,” and “general ignorance and poverty” must have prevented very many from seeking redress against their wealthy and powerful masters.\nobreak\textsuperscript{36} Despite such unfavorable odds, alleged slaves began claiming property rights to their homesteads in that very year. The boldness and novelty of such claims were noted by missionaries who had just been allowed into British territories by the East India Charter Act of 1813.\nobreak\textsuperscript{37}\par{\small 33. An early and influential account is Ranajit Guha, \textit{Elementary Aspects of Peasant Insurgency in Colonial India} (Oxford University Press, 1983). Recent critiques of Guha’s approach for the period that this essay focuses on include Jon E. Wilson, “‘A Thousand Countries To Go To’: Peasants and Rulers in Late Eighteenth-Century Bengal,” \textit{Past and Present} 189 (2005): 81–109; and Andrew Sartori, \textit{Liberalism in Empire: An Alternative History} (University of California Press, 2014).\par}{\small 34. Amal Kumar Chattopadhyay, \textit{Slavery in the Bengal Presidency, 1772–1843} (Golden Eagle Publishing House, 1977); Chatterjee, \textit{Gender, Slavery, and Law in Colonial India}; Major, \textit{Slavery, Abolitionism, and Empire in India.}\par}{\small 35. Chattopadhyay, \textit{Slavery in the Bengal Presidency}, 44.\par}{\small 36. Quoted in ibid., 44. A similar argument about the difficulty of rebellion in early colonial India was made by Ranajit Guha, “The Prose of Counter-Insurgency,” in Ranajit Guha, ed., \textit{Subaltern Studies II: Writings on South Asian History and Society} (Oxford University Press, 1983), 1–42.\par}{\small 37. William Adam, \textit{The Law and Custom of Slavery in British India, in a Series of Letters to Thomas Fowell Buxton, Esq.} (Weeks, Jordan, \& Co., 1840), 17–20.\par}
{\scriptsize\color{gray} [p. 9]}\ The judgment passed in the conflict discussed above became a precedent. In another instance in Dhaka in 1831, a slaveholder—Gauri Sankar Roy—filed a suit to recover runaway slaves.\nobreak\textsuperscript{38} Two of the defendants—Kirti Narayan and Sri Narayan Deo—appeared to testify and made identical claims about being tenants rather than slaves. The Narayan brothers insisted that their father, Binod Ram, was “free and earned his livelihood as a cultivator and tenant of land and by service.”\nobreak\textsuperscript{39} After the demise of their father, the brothers continued to hold lands as tenants on the estate of the plaintiff and paid rent. Notably, the plaintiff in his reply alleged that by the custom of the country, if \textit{Bhandaris} or household slaves cultivated for rent lands of their master more than those assigned for their support, such fact did not discharge them of their liabilities as slaves. On 24 August 1832, the Principal Sadar Amin agreed with the plaintiff and ordered all thirteen defendants to be made over to Gauri Sankar as his slaves. Subsequently, Kirti Narayan went in person to the Sadar Dewani Adalat on 21 February 1833 and asked for a special appeal to be admitted. Crucially, the ground for his request was that the decision in this case was contrary to the judgment passed by the SDA on 5 May 1832 in the previous case, that is, Golak Narayan Ray vs. Kewal Ram Deo and others. The SDA’s hand was forced, and the appeal was admitted by Mr. Rattray on 23 March 1833. Gauri Sankar Roy—the slaveholder—was asked to justify his claim. In response, Roy asserted that the appellants had received not only lands from him for their sustenance but also rations and other aids. He also insisted that the “want of a bill of sale after the lapse of ages” could not affect his right as a master. Slavery of the appellants was proved by the “continuity of hereditary service,” as well as the evidence adduced regarding the support received by the servants.\par{\small 38. Indian Law Commission, \textit{Report on Slavery, Volume 2}, 261–65.\par}{\small 39. The colonial state consistently used the sociological category of “cultivator” interchangeably with the legal category of “tenant.” In practice, though, the two statuses did not always overlap, because as the revenue chain expanded, a legally defined tenant was often a rent-collector from under-tenants who cultivated the land. In the districts of eastern Bengal in the 1830s, this distinction is likely to have been less consequential than what it became by the early twentieth century. For a detailed narrative about how this unfolded and its political implications, see Sartori, \textit{Liberalism in Empire}, chapter 5.\par}
{\scriptsize\color{gray} [p. 10]}\ Mr. Rattray, however, was not swayed. He proposed to reverse the decree of the lower court and passed his judgment on 7 August 1833, pointing out that the lack of a deed of sale and the lack of proof that the appellants held lands in return for hereditary service made Roy’s claim weak. More importantly, the judge made a new argument by insisting that even if such lands had been held in return for service, once the slaves quit such lands, “they cease[d] to be liable to any claim of servitude.” Acknowledging that the appellants were probably \textit{Bhakta-Dasa} slaves, that is, those who offered themselves in servitude for the sake of food, Rattray nevertheless insisted that “when the slave for his food abandons the service, he becomes free.” These opinions were confirmed by Mr. G. Stockwell when he heard the case on 7 December 1835, and thus the alleged slaves were emancipated.\par
{\scriptsize\color{gray} [p. 10]}\ An almost identical case reached the SDA in 1835 from the frontier district of Mymensingh, where a monumental peasant rebellion (1826–1833) had just concluded.\nobreak\textsuperscript{40} The slaveholder claimed that the defendants had served him as slaves in return for access to land. The alleged slaves rejected such an ascription and claimed that they had occasionally served the appellant and used some land, but only “in place of wages.” Now, after some years of service, they had decided to move to another village because they were treated there as “tenants.” A familiar back-and-forth between several courts and judges followed. Eventually, the case arrived at the SDA, and the final decision was passed by Mr. Stockwell on 17 May 1836. He acknowledged that the defendants may have been slaves in consideration of lodging and lands for support. It was also clear, however, that they had “abandoned such lodging and support.” Citing his judgment of 7 December 1835 in the case of Gauri Sankar Roy vs. the Narayan brothers as a precedent, Stockwell concluded that the respondents were “slaves of the class of slaves for their food.” On surrender of the lands held, therefore, they were “entitled to emancipation.” Soon after, another identical case from Sylhet was admitted into the SDA, although the decision in that instance did not go in favor of the alleged slaves for technical reasons.\nobreak\textsuperscript{41}\par{\small 40. Indian Law Commission, \textit{Report on Slavery, Volume 2}, 269–71. On the Mymensingh peasant rebellion, see Willem van Schendel, “‘Madmen’ of Mymensingh: Peasant Resistance and the Colonial Process in Eastern India, 1824–1833,” \textit{The Indian Economic and Social History Review} 22, 2 (1985): 139–73; and Gautam Bhadra, \textit{Iman O Nishan: Unish Shatake Banglar Krishak Chaitanyer ek Adhyay}, \textit{c.1800–1850} (Center for Studies in Social Sciences, 1994), 20–164.\par}{\small 41. Indian Law Commission, \textit{Report on Slavery, Volume 2}, 276–80.\par}
{\scriptsize\color{gray} [p. 10]}\ The trajectory of these cases shows that there was no consensus among judges and hence no coherent colonial policy regarding legal disputes over the status of “slaves.” Indeed, the Indian Law Commission’s report acknowledged that there was a great variety of opinion and practice on this matter and such variety was observable not only in different and distant parts of the country, but even in neighboring districts, and “sometimes [in] different divisions of the same district, and even at different times in the same court.”\nobreak\textsuperscript{42} Some officials admitted the legal standing of slavery and the rights and obligations arising out of it, whereas others rejected it entirely. And even among those who recognized slavery, “some would confine its existence within the bounds prescribed by the Mahomedan and Hindoo laws, others would extend it over the wider field of custom and usage; as Magistrates, some would render assistance to the master to recover his fugitive slaves, others would abstain from any interference; some would permit to the master the right of moderate correction, others would not countenance the exercise of such authority.”\nobreak\textsuperscript{43} The law, in short, was in a “defective” state, and it lacked “clear rules” that could guide judgments.\par{\small 42. Indian Law Commission, \textit{Report on Slavery, Volume 1}, 117.\par}{\small 43. Ibid., 118.\par}
{\scriptsize\color{gray} [p. 11]}\ Moreover, the defective state of the law and the ensuing confusion were partly a function of contradictions \textit{internal} to “Hindu” law. Based on classical texts including those of Manu, Vrihaspati, and Narada, Brahman (the uppermost, priestly caste) \textit{pandits} (scholars) had provided the colonial state with a list of fifteen classes of legal slaves in Hindu law.\nobreak\textsuperscript{44} That list, however, raised more questions than it answered. For example, how would one distinguish in practice between \textit{Bhakta-Dasa} slaves, or those who offered themselves in servitude for the sake of food, and \textit{Atma-Bikrita} slaves, that is, those who sold themselves for a price? People who entered servitude out of economic distress often received more than just food from their masters, and the question remained whether those additional benefits should be considered a “price” or not. The distinction mattered because, according to the pandits themselves, a \textit{Bhakta-Dasa} slave could be emancipated upon relinquishing his subsistence, but a self-sold slave could never become free without the will of the master.\nobreak\textsuperscript{45} In the cases narrated above, the masters did not conceive of their slaves as \textit{Bhakta-Dasas}, and neither did some colonial officials, but fortunately for the alleged slaves, the highest-ranking judges did.\par{\small 44. For the full list, see Chattopadhyay, \textit{Slavery in the Bengal Presidency}, 177–78. Cf. Indian Law Commission, \textit{Report on Slavery, Volume 2}, 268–69.\par}{\small 45. See ibid., 268; and Adam, \textit{The Law and Custom of Slavery in British India}, 17–20.\par}
{\scriptsize\color{gray} [p. 11]}\ Such problems of definition could easily become more complicated. How, for instance, would one distinguish between either \textit{Bhakta-Dasa} or \textit{Atma-Bikrita} slaves and \textit{Rina-Dasa} slaves, or distressed debtors who voluntarily engaged themselves to their creditors for a stipulated period? All three categories or a combination of them could easily be applicable in a particular case, but the rules of emancipation (according to the Brahmans) were different for the different categories. A remark on this point by the Officiating Magistrate of Rangpur is insightful: “In ninety cases out of the hundred, the servile status of slaves originated either by their forefathers having made themselves bondsmen by borrowing small sums of money (which bondage ought to have expired with their lives), or the descendants of cultivators who have died in debt, either to the zamindars, or persons holding small farms under them: as experience has proved to me, neither would scruple to compel the widows or children of their deceased bondsmen, or insolvent deceased cultivators, to give them written engagements of being slaves or bondsmen for life; and this instrument would cause the parties and their descendants to become slaves in perpetuity.”\nobreak\textsuperscript{46}\par{\small 46. Indian Law Commission, \textit{Report on Slavery, Volume 1}, 66.\par}
{\scriptsize\color{gray} [p. 12]}\ It was thus difficult, if not impossible, to know for certain exactly how servile statuses originated, and whether they involved coercion or not. In many parts of Bengal, these kinds of cases were in the clear majority, which shows that the alleged slaves in the conflicts discussed above may well have been tenants at some point.\par
{\scriptsize\color{gray} [p. 12]}\ The trickiest legal question, however, was that of whether servile status (whatever its source) should be inherited by descendants. On this issue, at least, the Brahman pandits provided the colonial state with a clear guideline: a \textit{Griha-Jata}, or one born in the house of a master from a female slave, was also a slave, and \textit{could not be emancipated without the will of the master.} The definition was extremely wide because the child did not need to be of the master to be considered a slave, nor did the type of slavery ascribed to the child’s father matter. Moreover, women cohabiting with any type of slave necessarily acquired servile status. Without doubt, this was a strategy to reproduce servitude, but it is clear that by the 1820s and 1830s, descendants of erstwhile slaves did not want to put up with it. In all of the cases discussed above, the masters were trying to ascribe servile status to the descendants of those who \textit{may} have been slaves to their families, but the new generations resisted. When colonial judges decided to emancipate the alleged slaves in these cases by treating them as \textit{Bhakta-Dasas} rather than \textit{Griha-Jata}, they were exceeding the space of “Hindu” law, even if they did not always acknowledge it.\par
{\scriptsize\color{gray} [p. 12]}\ In other cases, too, this problem reared its head, and the decisions passed by the colonial state were significant. In 1831, the courts had to resolve a conflict from a district that falls in modern-day Jharkhand and borders western Bengal. The plaintiff alleged that his forebears had purchased some slaves during the great famine of 1769–70 and that he was filing the suit to recover as his slaves those who were the descendants of the original \textit{Atma-Bikrita} slaves two generations removed.\nobreak\textsuperscript{47} In response, the defendants claimed that their ancestors had never sold themselves into slavery but rather had received “loans” of coarse grains and other foodstuffs that had since been repaid. The plaintiff shot back, saying that in a famine “no one supported another, \textit{particularly one of low caste}, for the mere acknowledgment of grain or money.”\nobreak\textsuperscript{48} Eventually, the case was decided against the plaintiff because the bill of sale provided was written on plain paper, which the courts found suspicious.\par{\small 47. Indian Law Commission, \textit{Report on Slavery, Volume 2}, 258–60.\par}{\small 48. Ibid., emphasis mine.\par}
{\scriptsize\color{gray} [p. 12]}\ Another conflict from the frontier district of Sylhet in 1826 had an identical structure: the slaveholder alleged desertion, but the alleged slaves responded that they were tenants.\nobreak\textsuperscript{49} In this case, a genuine bill of sale from 1757 was produced. The question was whether the presence of a clause stating that the person was selling himself for “sixty and seventy years of service” meant that after that term, his remaining descendants should be considered free. One judge claimed that the limitation of time was merely in conformance with “custom” and was meant to ensure the “exemption from labor in old age, not freedom.” Others countered that when a person sold himself for a stipulated period, it could scarcely imply that his wife, son, and further descendants “should pass as owned persons, and be liable to render service to the issue of the hirers.” The alleged slaves were eventually emancipated, but this was done \textit{despite} the suggestions of the Brahman pandits, expressly regarding this particular case, that the alleged slaves were \textit{Griha-Jata} and \textit{could not be emancipated without the will of their master.} \nobreak\textsuperscript{50}\par{\small 49. Ibid., 265–68.\par}{\small 50. Ibid., 268, emphasis mine.\par}
{\scriptsize\color{gray} [p. 13]}\ We see, then, that at least some colonial officials were upholding the undeniably modern idea that individuals cannot be held responsible for the sins and debts of their forebears. When judges decided to emancipate slaves, they usually did so either by playing on the nebulousness and overlapping of definitions or by going beyond the ambit of Hindu law as prescribed by the Brahmans. It is also clear, however, that this was not colonial benevolence but rather a response to pressures from below by alleged slaves. The evidence suggests that descendants of erstwhile slaves primarily sought property rights over their bodies and labor. They wanted to be able to choose where, how, and under whom they would work, and they preferred contractual relationships of “loans” and “wages” over the paternalism of a master-slave relationship. This was a rejection of caste norms as well because we know that at least in the far-eastern districts of Mymensingh and Sylhet, nankar lands were given for subsistence to the lowest castes.\nobreak\textsuperscript{51} The disputes show that the alleged slaves wanted to be able to move freely whenever they wanted to, especially in eastern Bengal, where land was still plentiful, and where the frontier was closed only in the late nineteenth century.\nobreak\textsuperscript{52} Slaves, in short, preferred the legal status of tenants or “raiyats,” even if that meant being \textit{paikasht} raiyats for a time, that is, raiyats who migrate often and are not permanently settled in a particular village.\nobreak\textsuperscript{53}\par{\small 51. Bhadra, \textit{Iman O Nishan}, 24–25.\par}{\small 52. For an overview of Bengal’s agrarian history under colonial rule, see Sugata Bose, \textit{Peasant Labor and Colonial Capital: Rural Bengal Since 1770} (Cambridge University Press, 1993).\par}{\small 53. Although I have focused on the actions of upper-caste, Hindu landowners, Muslim landowners also tried to prevent their “slaves” from migrating. For an example, see Indian Law Commission, \textit{Report on Slavery, Volume 2}, 272–73. The colonial state found it easier to apply Islamic law to cases of slavery because according to the Maulvis, the only legal classes of slaves were those who were captured in war—either a general war or a holy war against “infidels.”\par}
\Needspace{10\baselineskip}\section*{Making Sense of the Evidence: Abolition and Counter-Revolution in Colonial Bengal, 1772–1843}
{\scriptsize\color{gray} [p. 13]}\ To interpret this evidence, we must consider the trajectory and contradictions of colonial governance, as well as the \textit{longue durée} history of the society being governed. Certainly, ambivalence about questions of slavery and caste was a constant feature of colonial policy from the 1770s onward. There was an abolitionist impulse, which found expression in periodic regulations aimed at curtailing the slave trade in the Indian Ocean. In 1774, commerce in slaves within the territories of the British East India Company (EIC) was banned. In 1789, the export of slaves was prohibited and imports were eventually outlawed in 1811.\nobreak\textsuperscript{54} Such measures were only partially successful because the slave trade was extremely lucrative, and the company failed to prevent its servants from participating in the trade and from owning slaves. The deeper problem, however, was the company’s legislative measures.\par{\small 54. For more details on these regulations and their effects, see Richard B. Allen, “Suppressing a Nefarious Traffic: Britain and the Abolition of Slave Trading in India and the Western Indian Ocean, 1770–1830,” \textit{The William and Mary Quarterly} 66, 4 (2009): 873–94; and Chattopadhyay, \textit{Slavery in the Bengal Presidency}, 79–115.\par}
{\scriptsize\color{gray} [p. 14]}\ In 1772, Governor-General Warren Hastings sanctioned slavery in British India as a penal measure.\nobreak\textsuperscript{55} The ostensible aim was to prevent “hereditary” robbers by enslaving entire families and putting them to work. As the Court of Directors put it, the robbers causing havoc across the company territories were “wretches who have placed themselves in a state of declared war with Government, and are, therefore, wholly excluded from every benefit of its laws.”\nobreak\textsuperscript{56} Crucially, when justifying the use of slavery to combat robbery, the state introduced the distinction between the “evil” character of Atlantic slavery and the “mild” character of its Indian counterpart. In India, apparently, slaves were “treated as the children of families to which they belong; and often acquire a \textit{much happier state by their slavery than they could have hoped for by the enjoyment of liberty.}”\nobreak\textsuperscript{57} Another mode of justification was the grounding of slavery in the “religious” laws of “Hindus” and “Muslims,” which necessarily implicated “caste” in the colonial discourse on (Hindu) slavery.\par{\small 55. Chattopadhyay, \textit{Slavery in the Bengal Presidency}, 10.\par}{\small 56. West Bengal State Archives (WBSA), Proceedings of the Committee of Circuit at Kasimbazar, Vol. I to III, 28 June 1772, 22–23.\par}{\small 57. Ibid., emphasis mine.\par}
{\scriptsize\color{gray} [p. 14]}\ The state first connected caste to religion in 1772, when it decided that in all matters pertaining to inheritance, marriage, and other religious usages and institutions, the laws of the Koran concerning Muslims and those of the “Shaster” (sic) with respect to the Hindus would be followed. The exposition of the law was left to the Maulvis for Muslims, and the Brahmans for Hindus, thereby entrenching an upper-caste point of view as law. In 1793, Hindu and Muslim laws of marriage, caste, succession, and inheritance were upheld again by the Marquis Cornwallis. But the question of whether Hindu and Muslim laws of \textit{slavery} should also be upheld, and whether slavery should be deemed a domain of “religious” practice of the two communities, was still open.\nobreak\textsuperscript{58} The matter was settled arbitrarily in 1798 when the Sadar Dewani Adalat decreed that the \textit{spirit}, if not the letter, of the 1772 and 1793 decisions suggested that Islamic and Hindu laws should also apply to cases of slavery, and this was confirmed by the Governor-General on 12 April 1798. Until legal non-recognition of slavery by Act V of 1843, this remained the shaky foundation upon which slavery in British India rested.\nobreak\textsuperscript{59}\par{\small 58. Adam, \textit{The Law and Custom of Slavery in British India}, 24–28.\par}{\small 59. Ibid., 26.\par}
{\scriptsize\color{gray} [p. 14]}\ In truth, the 1798 decision followed logically from the ones made in 1772 and 1793 because caste \textit{was} correlated to slavery, and if the former could be deemed a religious usage, then why not the latter? At the time, caste and slavery functioned as overlapping yet distinct modalities of labor control in Bengal: the number of “impure” castes set the upper limit to the total number of agrarian slaves, because “pure” castes could be domestic slaves to castes of the same or higher status, but they could not be agrarian slaves.\nobreak\textsuperscript{60} Witnesses belonging to multiple castes arrived at the same conclusion in their testimonies to the colonial state, and Brahman witnesses went to great lengths to emphasize that only some castes were “fit” for some kinds of work. Vishnu Dutt Dalli, for instance, testified on 5 February 1839 that slaves of the “low-ranked” Chandala, Dom, Hera, and Kumar castes could only be employed in outdoor work.\nobreak\textsuperscript{61}\par{\small 60. This was true in southern India as well. See Dharma Kumar, “Caste and Landlessness in South India,” \textit{Comparative Studies in Society and History} 4, 3 (1962): 337–63; and Viswanath, \textit{The Pariah Problem.}\par}{\small 61. Indian Law Commission, \textit{Report on Slavery, Volume 2}, 23–29.\par}
{\scriptsize\color{gray} [p. 15]}\ Within the framework of early colonial law, therefore, one could insist that slavery was \textit{not} a domain of religious practice only in two ways: either accept that caste was also not just a religious institution but one that had a fundamentally secular basis, or artificially separate the question of slavery from the larger set of social relations within which it was embedded. The former was never done, and the latter course was followed in the run-up to 1843.\par
{\scriptsize\color{gray} [p. 15]}\ Legal conflicts between alleged slaves and their masters became a pressing issue by the 1810s. John Richardson, the District Judge and Magistrate of Bundelkhand, and William Leycester, a judge of the Committee of Circuit, argued for complete abolition in 1808 and 1815, respectively.\nobreak\textsuperscript{62} Chief Justice John Harrington sharpened the case for intervention with a poignant question in 1818. How, he asked, could slaves hope to attain the liberty they desired when they were “possessed of and [could] acquire no property to enable them to institute or defend a suit” and when they could be subjected to restraint or even confinement while suits were pending?\nobreak\textsuperscript{63} The case for abolition, however, was summarily rejected by a state deeply invested in extracting revenue and maintaining “order.” Acting Governor-General John Adam (January–August 1823) and his successor William Amherst (1823–1828) both concurred that slavery was interwoven with the “customs and feelings” of the people. Intervention, therefore, would only wreak discontent and it was inexpedient to hazard the inconvenience.\nobreak\textsuperscript{64}\par{\small 62. Chattopadhyay, \textit{Slavery in the Bengal Presidency}, 171, 191.\par}{\small 63. Quoted in ibid., 194.\par}{\small 64. Ibid., 199–200.\par}
{\scriptsize\color{gray} [p. 15]}\ In the 1833 East India Charter Act, Parliament inserted a clause stating that slavery would be abolished in the Indian territories by 12 April 1837.\nobreak\textsuperscript{65} The EIC’s Court of Directors responded that they would only approve of legislation against slavery if it did not do “violence to the feelings of Caste, or the rights of property.”\nobreak\textsuperscript{66} The original clause was even modified to state that no regulation would be passed without “due regard to the laws of marriage, and the rights and authorities of fathers.”\nobreak\textsuperscript{67} This is not surprising because the terms of discussion had already been set when caste was relegated to the religious domain in 1772. This meant, paradoxically and perversely, that any acknowledgment of the caste-slavery relationship could only serve to strengthen the argument \textit{against} intervention and abolition.\par{\small 65. \textit{Slavery and the Slave Trade in British India; With Notices of the Existence of these Evils in the Islands of Ceylon, Malacca, and Penang, Drawn from Official Documents} (Thomas Ward and Co., 1841), 47–48.\par}{\small 66. Ibid., 47.\par}{\small 67. Ibid., 48–49.\par}
{\scriptsize\color{gray} [p. 15]}\ The Law Commission sanctioned by the 1833 Charter was set up in 1835, but the Report on Slavery—which is the chief evidentiary basis of this essay—was published after significant delay in 1841. Even then, the defenders of slavery in both India and Britain kept insisting that abolition would make landowners and noblemen indignant; that laborers would become idle and unproductive without “correction” and coercion; and that slavery in India was “nothing more than an affair of caste.”\nobreak\textsuperscript{68} Faced with such vicious opposition, the opponents of slavery chose the easy way out. Instead of arguing that relegating caste to the domain of religion had been a mistake from the beginning, abolitionists conceded the artificial distinction between caste and slavery. They ended up being narrowly legalistic, and they wrongly assumed that caste distinctions would cease to matter in a regime of “free labor.”\nobreak\textsuperscript{69} Act V, passed on 7 April 1843, argued only for the legal non-recognition of slavery and did not actively declare anyone free.\nobreak\textsuperscript{70} Even this half-hearted measure was applied selectively, and the colonial state continued to return alleged slaves to their masters. The principal aim, always, was the extraction of revenue and the control of labor.\par{\small 68. Ibid., 40–52, 49. Cf. Chattopadhyay, \textit{Slavery in the Bengal Presidency}, 156–251; and Indian Law Commission, \textit{Report on Slavery, Volume 1}, 371–79.\par}{\small 69. The Baptist and later Unitarian missionary William Adam was especially prone to this error. See Adam, \textit{The Law and Custom of Slavery in British India}, 172.\par}{\small 70. For the full text of Act V, see Chattopadhyay, \textit{Slavery in the Bengal Presidency}, 250–51.\par}
{\scriptsize\color{gray} [p. 16]}\ These trends show that the tendency to consign caste to the religious sphere and to collude with landed elites was not new in the late nineteenth century. The policy of “religious neutrality” after 1858 meant that caste was relegated entirely to a “national social” domain beyond politics,\nobreak\textsuperscript{71} but the impulse itself had deeper roots. The early colonial state \textit{created} precedents under “Hindu” law from scratch, whereas slavery in Mughal Bengal—especially in the eastern, frontier districts—was grounded in economic distress and was not reducible to textual dicta.\nobreak\textsuperscript{72} The debates between conservationists and abolitionists also show that the Orientalist trope of caste as a static, millennia-old “religious” institution expressing a uniquely Indian penchant for hierarchy was strong in early colonial India.\nobreak\textsuperscript{73} This view influenced Act V and it survived well into the twentieth century with negative consequences for Dalits.\par{\small 71. Viswanath, \textit{The Pariah Problem.}\par}{\small 72. Adam, \textit{The Law and Custom of Slavery in British India}, 54.\par}{\small 73. Nicholas B. Dirks, \textit{Castes of Mind: Colonialism and the Making of Modern India} (Princeton University Press, 2001). A twentieth-century re-statement of this supposed penchant is Louis Dumont, \textit{Homo Hierarchicus: The Caste System and its Implications}, trans. Mark Sainsbury, Louis Dumont, and Basia Gulati (Oxford University Press, 1970).\par}
{\scriptsize\color{gray} [p. 16]}\ This does not imply, however, that caste was merely a “colonial construction.”\nobreak\textsuperscript{74} Such a formulation occludes from vision longer-running South Asian social histories and the ethical-political valence of abolition in a context shaped by caste and agrarian slavery. It erases the ideological content of abolition \textit{tout court} and cannot explain the subjectivity of the alleged slaves. Methodologically, the colonial construction thesis presupposes that any acknowledgment of the \textit{longue durée} history of caste \textit{necessarily} implies an acceptance of the Orientalist trope. If, however, caste is understood in terms of differential access to property and power—instantiated in practice as differential access to land, water, and markets, and the regulation of sexuality—then a history of caste can be written that refuses the binary of \textit{either} millennia-old religious institution \textit{or} merely a product of colonial governance.\par{\small 74. Dirks, \textit{Castes of Mind.}\par}
{\scriptsize\color{gray} [p. 16]}\ There is a long and extensive tradition of historiography that approaches caste in such terms. It shows that in the early historical era (ca. 1000 BCE–ca. 300 CE) and the early medieval period (ca. 600–1200 CE), caste formation was congruent with agrarian expansion, the need to streamline \textit{inheritance} over the control of land, and the extraction of revenue from land and markets. In other words, the growth of stratification was coeval with the emergence of patriarchy and the formation of states with fiscal bases in agriculture.\nobreak\textsuperscript{75} Caste was thus a way to \textit{directly} control resources and \textit{indirectly} control labor. Kinship relations continued to mediate social relations of labor, which meant that chattel slavery was never of major importance in India, and belonging to a “lower” caste did not automatically mean slave status.\nobreak\textsuperscript{76} At the same time, some groups existed only on the fringes of the absorptive process of sedentarization. They were simultaneously excluded (both spatially and socially) and yet included within the system as super-exploited laborers. These “outcaste” or “untouchable” communities varied widely across South Asia in terms of their ethnic and cultural identities, and the concept of untouchability ultimately depends on a vague, Brahmanical notion of ritual pollution. Nevertheless, all Indian vernaculars recognize the existence of a segregated sub-population forced to perform degrading labor deemed “impure” by society at large.\nobreak\textsuperscript{77} Pan-Indian terms such as \textit{Chandala} do identical conceptual work.\par{\small 75. Influential contributions include Damodar Dharmananda Kosambi, \textit{An Introduction to the Study of Indian History} (Popular Prakashan, 1956); Ram Saran Sharma, \textit{Sudras in Ancient India: A Social History of the Lower Order Down to Circa AD 600} (Motilal Banarsidass, 1980); Suvira Jaiswal, \textit{The Making of Brahmanic Hegemony: Studies in Caste, Gender, and Vaishnava Theology} (Tulika Books, 2016); Romila Thapar, \textit{Ancient Indian Social History: Some Interpretations} (Orient Longman, 1978); Manu V. Devadevan, \textit{The ‘Early Medieval’ Origins of India} (Cambridge University Press, 2020); and Sumit Guha, \textit{Beyond Caste: Identity and Power in South Asia, Past and Present} (Brill, 2013).\par}{\small 76. Damodar Dharmananda Kosambi, “The Basis of Ancient Indian History (I),” \textit{Journal of the American Oriental Society} 75, 1 (1955): 35–45, 44.\par}{\small 77. See Rupa Viswanath’s contribution to “A Roundtable on Rupa Viswanath’s \textit{The Pariah Problem: Caste, Religion, and the Social in Modern India} and the Study of Caste,” \textit{Modern Asian Studies} 56, 1 (2022): 1–64, 49–64. A slightly dated but insightful discussion of competing theories about the origins of untouchability is Prabhati Mukherjee, \textit{Beyond the Four Varnas: The Untouchables in India} (Motilal Banarsidass, 1988).\par}
{\scriptsize\color{gray} [p. 17]}\ Early medieval Bengal was no exception to these trends. The two major ruling dynasties—the Palas (ca. 750–1162 CE) and the Senas (ca. 1070–1230 CE)—were explicitly committed to the maintenance of a normative caste order.\nobreak\textsuperscript{78} They made land grants on a large scale to upper castes, and the collection of revenue from these lands, cultivated by labor that was ranked as “low caste,” was the fiscal basis of the state. From the late ninth century onward, grants began to mention the entire caste spectrum with varying degrees of access to land: Brahmans were at the top and the “untouchable” Chandalas (the “lowest of the mortals”) were at the bottom.\nobreak\textsuperscript{79} Sources from the early centuries of the second millennium CE show that the lowest castes involved in agrarian and other forms of menial labor—such as the Chamars, the Chandalas, and the Doms—were kept spatially segregated, and they lived outside the boundaries of political capitals and upper-caste rural settlements.\nobreak\textsuperscript{80} We have no way of estimating variations in the size of these communities over time, but that is a function of silences in the archive. Recent scholarship on slavery in medieval South India and Mughal South Asia has emphasized the challenges posed by such processes of “invisibilization.”\nobreak\textsuperscript{81}\par{\small 78. For details on the role of religion in the creation of an agrarian state and society, see Kunal Chakrabarti, \textit{Religious Process: The Purāṇas and the Making of a Regional Tradition} (Oxford University Press, 2001).\par}{\small 79. Ryosuke Furui, \textit{Land and Society in Early South Asia: Eastern India 400–1250 AD} (Routledge, 2020).\par}{\small 80. Rina Bhaduri, “Madhyayuge Banglay Nagarbinyasher Dhara,” in Aniruddha Ray and Ratnabali Chattopadhyay, eds., \textit{Madhyayuge Banglar Samaj O Sanskriti} (K. P. Bagchi \& Co., 1992), 31–59, 53.\par}{\small 81. Leslie C. Orr, “Slavery and Dependency in Medieval South India,” in Craig Perry \textit{et al.}, eds., \textit{The Cambridge World History of Slavery, Volume 2, AD 500–AD 1420} (Cambridge University Press, 2021), 313–33; Emma Kalb, “In and Out of Sight: Textual Traces of Slavery and the Enslaved in Mughal South Asia,” in Mònica Ginés-Blasi, ed., \textit{Intentional Invisibilization in Modern Asian History: Concealing and Self-Concealed Agents} (De Gruyter, 2025), 121–39.\par}
{\scriptsize\color{gray} [p. 17]}\ Mughal India was a vastly different context and yet caste continued to be \textit{reconstituted} in new ways. The state was no longer explicitly committed to a normative caste order and the agrarian revenue system was sprawling and complex. At the village and sub-regional levels, however, a caste-based hierarchical structure remained the basic framework in much of the subcontinent.\nobreak\textsuperscript{82} Peasants had customary rights of usage, but they were tied to the land, which they could neither freely alienate nor desert at will.\nobreak\textsuperscript{83} In Bengal, the vast majority of field labor continued to be performed by middling and lower castes, including those deemed untouchable, whereas the largest estates were controlled by the ritually highest-ranked Brahmans and the powerful Kayasthas.\nobreak\textsuperscript{84} Indeed, as late as 1830, Brahmans continued to hold lands at 50\% lower rent than their lower-caste counterparts in some regions.\nobreak\textsuperscript{85}\par{\small 82. Irfan Habib, \textit{The Agrarian System of Mughal India, 1556–1707} (Asia Publishing House, 1963), 120–23, 136–89.\par}{\small 83. Ibid., 115–17.\par}{\small 84. Ratnalekha Ray, \textit{Change in Bengal Agrarian Society c.1760–1850} (Manohar Books, 1979), 30, 52–53, 89, 110–11, 136, 171, 210, 215–16, \textit{passim}; John R. McLane, \textit{Land and Local Kingship in Eighteenth-Century Bengal} (Cambridge University Press, 1993), 37, 54–55, 57, 62–63, 144, 156–57, 309–10, \textit{passim.}\par}{\small 85. Aditee Nag Chowdhury-Zilly, \textit{The Vagrant Peasant: Agrarian Distress and Desertion in Bengal, 1770–1830} (Franz Steiner, 1982), 144.\par}
{\scriptsize\color{gray} [p. 18]}\ Slavery in Mughal Bengal was reducible neither to caste nor to timeless “Hindu” and “Muslim” laws. Rather, it was grounded in economic distress caused by famines, floods, and debt.\nobreak\textsuperscript{86} Faced with a choice between servitude and death, many chose the former. Deeds of sale between 1694 and 1770 show that once a sale was effected, the master had complete control over a slave and their heirs. Punishment of runaway slaves was common. Masters aimed at the extraction of maximum possible work and reserved the right to cut loose unproductive slaves in old age.\nobreak\textsuperscript{87} Besides voluntary sale, the internal supply of slaves was also sustained through violence—that is, the kidnapping of men, women, and children—and natural growth, since women who married or cohabited with slaves also acquired servile status and their children inherited it.\nobreak\textsuperscript{88}\par{\small 86. Shamita Sarkar, “Slavery in Late Mughal Bengal,” \textit{Proceedings of the Indian History Congress} 55 (1994): 600–5; Habib, \textit{The Agrarian System}, 110.\par}{\small 87. Sarkar, “Slavery in Late Mughal Bengal,” 602.\par}{\small 88. Chattopadhyay, \textit{Slavery in the Bengal Presidency}, 12.\par}
{\scriptsize\color{gray} [p. 18]}\ At the same time, agrarian social relations were deeply embedded in caste norms, which shaped norms of slavery in turn. Not all castes could be asked to do all kinds of labor, and upper castes could not, in general, be agrarian slaves to those of lower status.\nobreak\textsuperscript{89} Prices varied by caste and plummeted during famines, and slaves belonging to menial castes such as the Bagdis had to explicitly disavow the rights of their children to any property that their parents had acquired.\nobreak\textsuperscript{90} Many deeds of sale were written in the vernacular, which shows that the norms and expectations of caste and slavery were fully legible to all actors within Indian society.\nobreak\textsuperscript{91}\par{\small 89. Indian Law Commission, \textit{Report on Slavery in India}, \textit{Volume 2}, 29, \textit{passim.}\par}{\small 90. Sarkar, “Slavery in Late Mughal Bengal,” 602; Chattopadhyay, \textit{Slavery in the Bengal Presidency}, 47.\par}{\small 91. Some vernacular bills of sale are compiled in Siva Ratan Mitra, \textit{Types of Early Bengali Prose} (University of Calcutta Press, 1922), 86–87, 105, 109, 115.\par}
{\scriptsize\color{gray} [p. 18]}\ Two apparently contradictory things could thus be simultaneously true. Economic, political, and cultural changes between the sixteenth and eighteenth centuries, which some scholars have associated with an “early modernity” in South Asia (or Eurasia more broadly), massively increased the stakes of being able to access land and markets more freely.\nobreak\textsuperscript{92} Commercialization and expanded division of labor meant that caste was in a state of flux and neither static nor merely religious.\nobreak\textsuperscript{93} Concurrently, elites attempted to tighten control over labor through slavery and a renewed emphasis on caste hierarchies. In Rajasthan and Maharashtra in western India, for example, the seventeenth and eighteenth centuries witnessed not only an increase in the incidence of slavery but also a renewed emphasis on “purity” of caste and “descent.”\nobreak\textsuperscript{94} The context in both provinces was one of regional assertions against Mughal rule and agrarian expansion. The slaves were not only war captives, but also those affected badly by famine and ongoing debt. They performed all kinds of labor, including military service and agrarian labor. Upward mobility among merchant castes also led to a desire for greater “distance” from those deemed untouchable, a project that was brought to fruition in alliance with the pre-colonial state of Marwar.\nobreak\textsuperscript{95}\par{\small 92. Sanjay Subrahmanyam, “Connected Histories: Notes Towards a Reconfiguration of Early Modern Eurasia,” \textit{Modern Asian Studies} 31, 3 (1997): 735–62; Rosalind O’Hanlon, “Contested Conjunctures: Brahman Communities and ‘Early Modernity’ in India,” \textit{The American Historical Review} 118, 3 (2013): 765–87. For a critique, see Dipesh Chakrabarty, “The Muddle of Modernity,” \textit{The American Historical Review} 116, 3 (2011): 663–75.\par}{\small 93. Sanjay Subrahmanyam, “Hearing Voices: Vignettes of Early Modernity in South Asia, 1400–1750,” \textit{Daedalus} 127, 3 (1998): 75–104. Cf. A. K. Ramanujan et al., eds., \textit{When God is a Customer: Telugu Courtesan Songs by Ksetrayya and Others} (University of California Press, 1994).\par}{\small 94. See the essays by Ramya Sreenivasan and Sumit Guha in Indrani Chatterjee and Richard M. Eaton, eds., \textit{Slavery and South Asian History} (Indiana University Press, 2006), 136–61, 162–86.\par}{\small 95. Divya Cherian, \textit{Merchants of Virtue: Hindus, Muslims, and Untouchables in Eighteenth-Century South Asia} (University of California Press, 2023). For a similar desire among merchants and middling agrarian groups in Bengal, see Hitesranjan Sanyal, “Continuities of Social Mobility in Traditional and Modern Society in India: Two Case Studies of Caste Mobility in Bengal,” \textit{The Journal of Asian Studies} 30, 2 (1971): 315–39.\par}
{\scriptsize\color{gray} [p. 19]}\ The early colonial state, by instituting a “despotism of law,” acted as a bulwark for this second trend against the first.\nobreak\textsuperscript{96} The conflicts over the legal status of alleged “slaves” must be read against this background. In the eastern, frontier districts of Bengal where conflicts recurred most frequently, caste formation was still inchoate. The Muslim peasantry that emerged in eastern Bengal from ca. 1574 onward was not composed of those who had earlier been “Hindu” but rather of autochthonous communities at the fringes of the agrarian order, including the Chandalas.\nobreak\textsuperscript{97} Upper-caste landed elites and moneylenders often forwarded the capital necessary for reclamation during the seventeenth and eighteenth centuries,\nobreak\textsuperscript{98} but the form that agrarian social relations could take was by no means a foregone conclusion. Slavery in this context was a tool for consolidating control.\nobreak\textsuperscript{99} Indeed, by the eighteenth century, large estates regularly included a considerable number of slaves.\nobreak\textsuperscript{100} In the eastern districts of Dhaka-Jalalpur, Rajshahi, and Sylhet, slaves were also sold \textit{with} the land.\nobreak\textsuperscript{101} Slaveholding, therefore, was especially important in frontier regions with a favorable land-labor ratio to guarantee labor supply. This is why, despite the extremely limited character of “abolition,” 580 Hindu and 700 Muslim landholders in Sylhet opposed Act V. In a petition submitted in February 1843, they claimed that hereditary slaves did not think of their situation as “degraded” and that men of “rank” had “no means of employing paid servants to perform the duties which were ordinarily performed by slaves.”\nobreak\textsuperscript{102}\par{\small 96. Radhika Singha, \textit{A Despotism of Law: Crime and Justice in Early Colonial India} (Oxford University Press, 1998).\par}{\small 97. Richard M. Eaton, \textit{The Rise of Islam and the Bengal Frontier, 1204–1760} (University of California Press, 1993), 118–19, 132.\par}{\small 98. Ibid., 220–27.\par}{\small 99. The data in Chattopadhyay, \textit{Slavery in the Bengal Presidency}, 31–78, shows not only the large number of slaves in early colonial Bengal but also the high proportion of agrarian slaves in the total number. By 1841, Bengal had approximately 2.5 million slaves, and the whole of British India likely had eight or nine million. Ibid., 66, 11.\par}{\small 100. Shirin Akhtar, “The Role of the Zamindars in Bengal, 1707–1772” (PhD diss., School of Oriental and African Studies, University of London, 1973), 249.\par}{\small 101. Indian Law Commission, \textit{Report on Slavery in India}, \textit{Volume 1}, 92.\par}{\small 102. Quoted in Chattopadhyay, \textit{Slavery in the Bengal Presidency}, 248–50.\par}
{\scriptsize\color{gray} [p. 20]}\ Thus, the alleged “slaves” in legal conflicts challenged not only slavery but also the \textit{practices} that produced and reproduced caste: the control of access to land and labor. The challenge was not articulated in language or texts. Nor does the evidence imply a coherent ideology of abolitionism. Nevertheless, action and thought cannot be artificially separated. By challenging established practices of slavery against all odds, the enslaved thereby \textit{effected} a critique of caste in practice. Arguably, such a practical challenge was more worrisome for those in power than a textual exegesis that did not necessarily challenge the practices that constituted caste. Why else, we might ask, were landholders in frontier districts so concerned about the toothless “abolition” of the colonial state?\par
{\scriptsize\color{gray} [p. 20]}\ Reading evidence from the colonial archive in terms of a larger structure of inference is not uncommon in a field of study where direct access to the mental life of historical subjects is virtually impossible. In his influential work, Thomas Holt also inferred what the emancipated black workers \textit{thought} by looking at what they \textit{did.} Although Holt had no evidence of actual statements made by the emancipated, he asked what kind of commitments could help us make sense of the actions they undertook. Fraught as this interpretive exercise is, it enabled Holt to situate the Jamaican case within broader debates about abolition, racism, and capitalism. This essay attempts something similar on a smaller scale for India, and unlike Holt, I do have evidence where the alleged “slaves” speak, albeit through the mediation of the colonial state.\par
{\scriptsize\color{gray} [p. 20]}\ Does the fact of this mediation make the archive unusable? Not necessarily. Scholars of the Atlantic world have been engaged in a lively debate for decades about how to access the subjectivity of the enslaved.\nobreak\textsuperscript{103} There is a broad consensus that legal conflicts recorded in colonial archives make social fault lines legible, and at certain moments in such conflicts even the enslaved could “speak.”\nobreak\textsuperscript{104} Despite the difference in context, the stakes involved are strikingly similar for South Asia: writing about and interpreting the actions of entire communities that have been deliberately “invisibilized” in the archive. Like their Atlantic counterparts, scholars of Indian Ocean slavery have recently noted that careful attention to legal conflicts and the indexing of court cases might enable an archivally grounded “global subaltern (social) history.”\nobreak\textsuperscript{105} This essay seeks to contribute to and deepen such conversations.\par{\small 103. For a recent, synthetic discussion of this literature, see Trevor Burnard, \textit{Writing the History of Global Slavery} (Cambridge University Press, 2023), 40–56.\par}{\small 104. Ibid., 47–51.\par}{\small 105. Alexander Geelen et al., “Between Markets and Chains: An Exploration of the Experiences, Mobility, and Control of Enslaved Persons in Eighteenth-Century South-West India,” in Alicia Schrikker and Nira Wickramasinghe, eds., \textit{Being a Slave: Histories and Legacies of European Slavery in the Indian Ocean} (Leiden University Press, 2020), 75–98, 78.\par}
{\scriptsize\color{gray} [p. 21]}\ At the same time, the colonial archive was politically constituted and historically specific. Although the Charter Act of 1813 has received far greater attention from historians due to its major implications for the future of Indian industry, the Act of 1833 was no less consequential politically.\nobreak\textsuperscript{106} It transferred the entire responsibility of governance and policymaking to a company-state invested in revenue extraction, and the India Question dropped out of the British Parliament completely. It also espoused an authoritarian vision of “reform” led by Thomas Macaulay, who chaired the Law Commission that produced the Report on Slavery. Thus, when Lord Brougham raised the issue of Indian slavery in Parliament on 5 October 1841, Lord Ellenborough responded that “Parliament had wisely directed that the initiation in legislation on this subject should be left with the Governor-General and council.”\nobreak\textsuperscript{107} Indeed, Ellenborough, who oversaw the passing of Act V as Governor-General of India between February 1842 and June 1844, had denounced the clause in the 1833 Act regarding abolition as “useless and unnecessary.” It would, he thought, “excite the indignation of the landed proprietors” and “shake the confidence and the allegiance of the native officers.”\nobreak\textsuperscript{108}\par{\small 106. Anthony Webster, “The Political Economy of Trade Liberalization: The East India Company Charter Act of 1813,” \textit{The Economic History Review} 43, 3 (1990): 404–19; Joshua Ehrlich, “The Crisis of Liberal Reform in India: Public Opinion, Pyrotechnics, and the Charter Act of 1833,” \textit{Modern Asian Studies} 52, 6 (2018): 2013–55.\par}{\small 107. House of Lords, \textit{Hansard}, 3rd ser., vol. 59 (5 October 1841), cols. 1115–28.\par}{\small 108. Quoted in \textit{Slavery and the Slave Trade in British India}, 49.\par}
{\scriptsize\color{gray} [p. 21]}\ The early colonial state, therefore, was a force of reaction, but such forces were ascendant at the time not only in India but throughout the wider imperial world \textit{and} within Europe and Britain. Most officials who governed India—even “Whigs” such as Macaulay who considered it necessary to appeal to the interests of the rising middle classes and “the democracy”—were political proponents of “order” and reform from above.\nobreak\textsuperscript{109} They opposed democratic, working-class movements everywhere including Chartism in Britain. Thus, Lord Ellenborough called the movement for electoral reform “a monstrous democracy,” and Robert Peel, who was Prime Minister (1841–1846) in 1843, called it “a reckless democracy.”\nobreak\textsuperscript{110} The reform that did occur secured the influence of the landed elite in Parliament and separated a section of the middle classes from proletarian workers by including the former within the political system.\nobreak\textsuperscript{111} Abolition, in this context, functioned like “moral capital” for a British Empire seeking atonement for the “sin” of slavery, which had led, many believed, to the loss of the American colonies.\nobreak\textsuperscript{112} Antislavery appealed to political elites because it enabled Parliament to keep flying the flag of liberty at a time when it was clamping down on liberties both at home and abroad.\nobreak\textsuperscript{113}\par{\small 109. Peter J. Gurney, “The Democratic Idiom: Languages of Democracy in the Chartist Movement,” \textit{The Journal of Modern History} 86, 3 (2014): 566–602.\par}{\small 110. Quoted in ibid., 571.\par}{\small 111. Hilton, \textit{A Mad, Bad, and Dangerous People?}, 420–36, 612–21.\par}{\small 112. Christopher L. Brown, \textit{Moral Capital: Foundations of British Abolitionism} (University of North Carolina Press, 2006).\par}{\small 113. Linda Colley, \textit{Britons: Forging the Nation, 1707–1837} (Yale University Press, 1992), 350–60.\par}
{\scriptsize\color{gray} [p. 21]}\ By thus historicizing the colonial archive, we can grasp the contradictions of abolition in South Asia in a dialectical manner. In an agrarian society shaped by caste and slavery, formal freedom and personal liberty mattered. Thus, the claim by Dalit studies scholars about the utility of “colonial” liberalism and modernity for Dalits is historically valid. In Bengal, Lockean claims about the constitution of property by labor shaped agrarian politics from the 1860s onward, well into the twentieth century.\nobreak\textsuperscript{114} Such claims were already present in embryonic form in the \textit{Paglai Dhum} peasant rebellion (1825–1833) in the eastern frontier district of Mymensingh. Here, hill tribes such as the Garos, the Hajongs, and the Hadis, who were being incorporated into a settled agrarian order as \textit{laboring} castes, resisted increased revenue assessment. They claimed that since it was through their labor (or of their forefathers) that forests had been reclaimed and the land made to bear fruit, the incremental profit and higher rent that resulted from the increased settlement in the region should accrue to them and not to the landlords or the state.\nobreak\textsuperscript{115} Indeed, the connection between land, labor, and property was first made by the proponents of Faraizism, an Islamic reformist movement that emerged to prominence among the weavers and cultivators of eastern Bengal in the 1820s.\nobreak\textsuperscript{116} Their slogan—\textit{langal jar, jami tar}, that is, the land belongs to the owner of the plough that tills it—inspired communist-led agrarian movements in Bengal throughout the twentieth century.\par{\small 114. Sartori, \textit{Liberalism in Empire.}\par}{\small 115. Bhadra, \textit{Iman O Nishan}, 20–164.\par}{\small 116. Sartori, \textit{Liberalism in Empire}, 111.\par}
{\scriptsize\color{gray} [p. 22]}\ By claiming that \textit{all} labor irrespective of caste status should be seen as constituting property equally, peasants articulated a profound critique of caste as a regime where different categories of persons were understood to be engaged in qualitatively distinct forms of labor, and those who labored the most were perpetually propertyless. Such claims, though, could only be made by tenants and not by slaves treated as property. Given that legal conflicts recurred most frequently in the eastern districts with surplus land, there can be little doubt that this was one of the reasons why slaves wanted to assert a property right over their own bodies and labor. The relevance of abolition in South Asia, therefore, cannot be reduced to the contingent requirements of colonial governance and empire.\nobreak\textsuperscript{117}\par{\small 117. The most comprehensive statement of the complicity between liberalism and empire in South Asia is Uday Singh Mehta, \textit{Liberalism and Empire: A Study of Nineteenth-Century British Political Thought} (University of Chicago Press, 1999). For a different view that foregrounds the role of metropolitan politics in shaping political thought, see Javed Majeed, \textit{Ungoverned Imaginings: James Mill’s ‘The History of British India’ and Orientalism} (Clarendon Press, 1992).\par}
{\scriptsize\color{gray} [p. 22]}\ Nevertheless, Act V, and liberalism in South Asia more broadly, functioned as tools of empire in a historical moment of counter-revolution that laid bare the emergent contradictions of capitalism. Massive proletarianization had brought the question of mass democracy and socialist transformation to the fore in Europe, which meant, as Marx insisted, that classical bourgeois revolutions were no longer feasible. There could only be counter-revolutions defending already constituted property, or social-republican revolutions against the capitalist order.\nobreak\textsuperscript{118}\par{\small 118. Marx, “The Bourgeoisie and the Counter-Revolution I–IV,” 210–12.\par}
{\scriptsize\color{gray} [p. 22]}\ This point brings to light the limits of abolitionists in South Asia who conceded the artificial distinction between caste and slavery. By the late eighteenth century, the practices that ensured differential access to property and power had created a \textit{de facto} separation between propertyless and propertied castes, which provided the grounds for the emergence of specifically capitalist labor relations. This can be discerned from the fact that in the early nineteenth century, a floating, rural proletariat existed in most parts of India, and, when industrialization began, the necessary labor was drawn primarily from this floating proletariat consisting of the lowest castes.\nobreak\textsuperscript{119} Therefore, scholars who argue that the lack of classic primitive accumulation in India is one reason for the distinctive character of capitalism in the region are only partly correct, because they ignore the long-running history of caste and the incorporation and reconstitution of caste within capitalism.\nobreak\textsuperscript{120} Lockean claims to property were ill-equipped to respond to these complexities, which explains why such arguments could function as a defense of the status quo—that is, a caste-differentiated distribution of property—by the 1910s.\nobreak\textsuperscript{121}\par{\small 119. Kumar, “Caste and Landlessness in South India”; Ranajit Dasgupta, “Factory Labor in Eastern India: Sources of Supply, 1855–1946; Some Preliminary Findings,” \textit{The Indian Economic and Social History Review} 13, 3 (1976): 277–329; Ian J. Kerr, \textit{Building the Railways of the Raj, 1850–1900} (Oxford University Press, 1995), 92.\par}{\small 120. See Kalyan Sanyal, \textit{Rethinking Capitalist Development: Primitive Accumulation, Governmentality, and Post-Colonial Capitalism} (Routledge, 2007). Cf. David Mosse, “The Modernity of Caste and the Market Economy,” \textit{Modern Asian Studies} 54, 4 (2020): 1225–71.\par}{\small 121. Rao, \textit{The Caste Question}, chapter 2.\par}
{\scriptsize\color{gray} [p. 23]}\ Though ultimately unsuccessful, the slogan of “forty acres and a mule” to the formerly enslaved was at least raised in the era of Reconstruction in the southern United States.\nobreak\textsuperscript{122} Something equally radical would have been necessary in India to resemble anything like a level playing field, but the counter-revolution of property, backed by an extractive colonial state, was even stronger. These contextual and political factors explain the failures and limits of liberty and abolition in South Asia, not the inability of the colonial state to understand the “slavery-to-kinship continuum,” or a purported attempt to smuggle in “European concepts of liberty and ‘possessive individualism.’”\nobreak\textsuperscript{123}\par{\small 122. Du Bois, \textit{Black Reconstruction}, 182–236, 580–636.\par}{\small 123. Chatterjee, \textit{Gender, Slavery, and Law in Colonial India}; Chatterjee, “Renewed and Connected Histories: Slavery and the Historiography of South Asia,” in \textit{Slavery and South Asian History}, 17–43, 33. Chatterjee leaves out agrarian servitude entirely from her analysis and focuses solely on domestic slaves of elite households. Her argument about a slavery-kinship continuum is based on an uncritical acceptance of the “absorptionist analysis” of African slavery by Igor Kopytoff and Suzanne Miers. For a review of this approach and a persuasive critique, see Frederick Cooper, “The Problem of Slavery in African Studies,” \textit{The Journal of African History} 20, 1 (1979): 103–125.\par}
\Needspace{10\baselineskip}\section*{Conclusion}
{\scriptsize\color{gray} [p. 23]}\ The contemporary literature on Indian Ocean slavery speaks to many themes broached in this essay. Scholars have shown that wages and loans, and the formal freedom they embodied, were not unimportant. Wages were used to purchase manumission in the Cape colonies, and freedom in the Indian Ocean often meant working for a wage.\nobreak\textsuperscript{124} “Flight” represented “fight,” and the lack of large-scale slave rebellions compared to the Atlantic does not imply a lack of resistance.\nobreak\textsuperscript{125} Rebellion against extractive empires and deeply rooted indigenous power structures was not easy, and any resistance we find in the archive requires a historical explanation.\par{\small 124. Kate Ekama, “Bound to be Free? Manumission in Cape Town, 1825–1834,” \textit{Slavery and Abolition} 46, 2 (2025): 441–62; Chakraborty, “Slave Trading and Slave Resistance in the Indian Ocean World.”\par}{\small 125. Rediker et al., \textit{A Global History of Runaways}; Matthias van Rossum and Jeannette Kamp, eds., \textit{Desertion in the Early Modern World: A Comparative History} (Bloomsbury Academic, 2016).\par}
{\scriptsize\color{gray} [p. 24]}\ New research on global histories of slavery, however, has not inspired attempts to write global histories of antislavery. The Atlantic-centered debate is rarely invoked today, and the dialectical approach of major contributors to that debate has been jettisoned in favor of a purely negative view of capitalism as \textit{necessarily} imperial and racist. The obverse of this tendency is the conceptualization of “resistance” as challenges only to extractive European empires, which leaves open the possibility of an uncritical affirmation of pre-colonial South Asian societies.\par
{\scriptsize\color{gray} [p. 24]}\ By contrast, this essay has argued that combining insights from the antislavery debate with the extensive literature on caste in India offers two benefits. One, it brings into sharp relief contradictions within Indian society, and thereby shows how and why abolition and the formal freedom of wage-labor mattered in that context. Those alleged to be “slaves” expressed dissatisfaction with the legal or \textit{de jure} status of complete ownership, whereas their indigenous masters argued that such rights existed \textit{de facto} in Bengali society. The “slaves” insisted that they had sold their \textit{capacity} to labor in return for wages and loans, and not their personhood. Such claims were motivated by a desire for personal liberty: a desire to \textit{become} something other than a “slave,” based on a recognition of the \textit{stakes} involved in it.\par
{\scriptsize\color{gray} [p. 24]}\ The colonial state and the global, capitalist context of counter-revolution smothered these aspirations. By thus historicizing both the potentials and failures of liberty, we can think about comparison in a grounded manner, one that does not proceed from \textit{a priori} assumptions either about the equivalence of “caste” and “race,” or about a \textit{necessary} connection between capitalism, slavery, and racism.\nobreak\textsuperscript{126} Close attention to the specificities of caste, slavery, and social aspirations within South Asia, and archivally grounded arguments about how and when they came to be mediated by the history of capitalism, can obviate the need to fit regional histories into yet another Atlantic-centric notion such as “racial capitalism.” At stake is a deeper understanding of how global capitalism has unevenly incorporated pre-capitalist social forms within it and reconstituted them over time.\par{\small 126. In a recent review essay, Sanjay Subrahmanyam has argued that hasty comparisons threaten to flatten out the historical specificities of both caste and racism. I take Subrahmanyam’s concerns seriously but disagree with his implicit suggestion that comparison may be impossible. See Sanjay Subrahmanyam, “Revisiting ‘Racial Capitalism,’” \textit{European Journal of Sociology} 64, 2 (2023): 173–81. Cf. n. 24 above.\par}
\Needspace{10\baselineskip}\section*{Acknowledgments}
{\scriptsize\color{gray} [p. 24]}\ I would like to thank three anonymous referees for their deep and critical engagement with the essay. Their suggestions led to substantial improvements in the structure and argument of the article. Andrew Sartori read several iterations of the piece over many years and offered helpful suggestions each time. Long-winded conversations with James M. Vaughn and Sunit Singh at The University of Chicago have influenced my thinking in more ways than I can acknowledge. All remaining errors are mine.\par
\end{document}
